\section{Medios y financiación}
Para ejecutar las tareas requeridas, el doctorando tiene accesos a una infraestructura tecnológica avanzada integrada en su entorno profesional, necesaria tanto para el procesamiento de trazas reales de producción como para el entrenamiento y validación de los modelos de ML propuestos, específicamente AWS ECS, S3 y RDS. Estos medios serán suficientes para implementar y probar los modelos predictivos propuestos para mejorar la eficiencia energética sin incurrir en costes añadidos por la adquisición de equipos físicos.
Adicionalmente, la investigación cuenta con acceso a las trazas reales de producción del Centre de Telecomunicacions i Tecnologies de la Informació (CTTI) de la Generalitat de Catalunya como caso de estudio principal, así como con conjuntos de datos públicos como Google Cluster Traces, lo que garantiza la disponibilidad de datos suficientes para entrenar y validar los modelos en condiciones reales.
La financiación se apoya en recursos propios.
El doctorando se encuentra en el cargo de Senior Analyst en proyectos estratégicos (Generalitat CTTI, Vodafone y Inditex).
Desde su posición dispone de una dedicación estable al proyecto de investigación y del soporte logístico para la validación requerida, sin necesidad de financiación externa.
